% ===== PRÉAMBULE CANONIQUE — à coller VERBATIM en tête de CHAQUE .tex =====
\documentclass[11pt,a4paper]{article}
\usepackage[utf8]{inputenc}\usepackage[T1]{fontenc}\usepackage{lmodern}
\usepackage[french]{babel}
\usepackage{amsmath,amssymb,mathtools}\usepackage{icomma}
\usepackage{siunitx}\sisetup{locale=FR}  % \qty{}{}, \si{}, \ang{}
\usepackage{xcolor}\usepackage{colortbl}
\usepackage{tikz}\usepackage{pgfplots}\pgfplotsset{compat=1.18}
\usepackage[siunitx,europeanresistors]{circuitikz} % circuits (élec.)
\usepackage[most]{tcolorbox}\usepackage{enumitem}\usepackage{array,booktabs}
\usepackage[a4paper,margin=2.2cm]{geometry}
\usetikzlibrary{arrows.meta,calc,angles,quotes,positioning,patterns,%
  backgrounds,shapes.geometric,decorations.pathmorphing,decorations.markings,intersections}
\definecolor{primary}{RGB}{222,49,99}\definecolor{primarydark}{RGB}{150,22,55}
\definecolor{defcol}{RGB}{0,114,178}\definecolor{methcol}{RGB}{52,92,148}
\definecolor{excol}{RGB}{0,135,90}\definecolor{attncol}{RGB}{200,40,55}
\tcbset{boxbase/.style={enhanced,breakable,boxrule=0.9pt,arc=2.5pt,
  left=3mm,right=3mm,top=2.3mm,bottom=2.3mm,fonttitle=\bfseries,coltitle=white}}
% --- boîtes TUTORIEL (noms EXACTS, ne pas renommer) ---
\newtcolorbox{cle}[1][]{boxbase,colback=defcol!6,colframe=defcol,
  title={\textsc{À retenir}\ifx\relax#1\relax\relax\else\textmd{ : \itshape #1}\fi}}
\newtcolorbox{methode}[1][]{boxbase,colback=methcol!7,colframe=methcol,sharp corners,
  title={\textsc{Méthode}\ifx\relax#1\relax\relax\else\textmd{ --- \itshape #1}\fi}}
\newtcolorbox{exemplebox}[1][]{boxbase,colback=excol!5,colframe=excol,
  title={\textsc{Exemple}\ifx\relax#1\relax\relax\else\textmd{ : \itshape #1}\fi}}
\newtcolorbox{piege}[1][]{boxbase,colback=attncol!6,colframe=attncol,
  title={\textsc{Piège}\ifx\relax#1\relax\relax\else\textmd{ : \itshape #1}\fi}}
\newtcolorbox{remarque}[1][]{boxbase,colback=black!4,colframe=black!45,
  title={\textsc{Remarque}\ifx\relax#1\relax\relax\else\textmd{ : \itshape #1}\fi}}
% --- boîtes EXERCICE / CORRIGÉ (noms EXACTS, auto-comptées) ---
\newtcolorbox[auto counter]{exo}[1][]{boxbase,colback=primary!4,colframe=primary,
  title={\textsc{Exercice}~\thetcbcounter\ifx\relax#1\relax\relax\else\textmd{ --- \itshape #1}\fi}}
\newtcolorbox[auto counter]{corrige}[1][]{boxbase,colback=excol!4,colframe=excol,
  title={\textsc{Corrigé}~\thetcbcounter\ifx\relax#1\relax\relax\else\textmd{ --- \itshape #1}\fi}}
\newcommand{\vv}[1]{\vec{#1}}\newcommand{\ds}{\displaystyle}
\newcommand{\R}{\mathbb{R}}\newcommand{\N}{\mathbb{N}}\newcommand{\Z}{\mathbb{Z}}
% (un \usepackage{fancyhdr}+en-tête est possible ; il est ignoré côté web)
% =========================================================================
\begin{document}
\sisetup{per-mode=symbol}

\begin{exo}[énergie potentielle de pesanteur]
Un livre de masse $m=\qty{2}{\kilogram}$ est posé sur une étagère à une hauteur $h=\qty{5}{\metre}$
au-dessus du sol (référence $E_p=0$ au sol). On prend $g=\qty{10}{\metre\per\second\squared}$.
\begin{center}
\begin{tikzpicture}[>=Stealth,scale=0.9]
  \draw[black!70,line width=0.8pt] (-0.3,0)--(4,0);
  \foreach \x in {-0.1,0.3,...,3.7} \draw[black!45] (\x,0)--(\x-0.18,-0.2);
  \draw[black!70,line width=1pt] (2.6,2.4)--(3.8,2.4);
  \draw[fill=primary!15,draw=black!70] (2.9,2.4) rectangle (3.4,2.75);
  \node at (3.15,2.57){\scriptsize $m$};
  \draw[<->,black!55,dashed] (2.0,0)--(2.0,2.4) node[midway,left]{\small $h=\qty{5}{\metre}$};
\end{tikzpicture}
\end{center}
Calcule l'énergie potentielle de pesanteur $E_p$ du livre.
\end{exo}

\begin{exo}[énergie potentielle élastique]
Un ressort de raideur $k=\qty{200}{\newton\per\metre}$ est comprimé d'une longueur
$x=\qty{0,10}{\metre}$. Calcule l'énergie potentielle élastique $E_p$ qu'il emmagasine.
\end{exo}

\begin{exo}[variation d'énergie potentielle]
On monte une caisse de masse $m=\qty{10}{\kilogram}$ de l'altitude $h_A=\qty{2}{\metre}$ à l'altitude
$h_B=\qty{8}{\metre}$ ($g=\qty{10}{\metre\per\second\squared}$).
\begin{enumerate}[label=\alph*)]
  \item Calcule la variation d'énergie potentielle $\Delta E_p=mg(h_B-h_A)$.
  \item Déduis-en le travail du poids $W(\vv G)=-\Delta E_p$. Est-il moteur ou résistant ?
\end{enumerate}
\end{exo}

\begin{exo}[loi de Hooke et énergie]
Pour un ressort, on mesure une force de rappel $F=\qty{6}{\newton}$ lorsqu'il est étiré de
$x=\qty{0,02}{\metre}$.
\begin{enumerate}[label=\alph*)]
  \item Calcule la constante de raideur $k$ à partir de $F=k\,x$.
  \item Calcule l'énergie potentielle élastique $E_p=\tfrac12 k x^2$ à cette élongation.
\end{enumerate}
\end{exo}

\begin{exo}[période et fréquence d'un ressort]
Un ressort de raideur $k=\qty{100}{\newton\per\metre}$ porte une masse $m=\qty{1}{\kilogram}$ et
oscille. On prend $\pi\approx3,14$.
\begin{enumerate}[label=\alph*)]
  \item Calcule la période $T=2\pi\sqrt{m/k}$.
  \item Déduis-en la fréquence $f=1/T$.
\end{enumerate}
\end{exo}
\end{document}
